\documentclass[10pt,a4paper]{amsart}
\usepackage[margin=19mm]{geometry}
\usepackage{amsmath,amssymb,mathtools}
\usepackage[scr=rsfs,cal=euler]{mathalfa}
\usepackage{xfrac}
\usepackage[unicode,hidelinks,bookmarks=false]{hyperref}
\newcommand{\ie}{i.e.\ }
\newcommand{\cf}{cf.\ }
\newcommand{\resp}{respectively\ }
\newcommand{\Subsection}{\subsection}
\providecommand{\nobreakdash}{\nobreak\hbox{-}}
\setlength{\emergencystretch}{2em}
\multlinegap0pt
\usepackage{bm}
\usepackage[shortlabels]{enumitem}
\usepackage[normalem]{ulem}
\newtheorem{proposition}{Proposition}
\newcommand{\R}{\mathbb{R}}
\newcommand{\lra}{\longrightarrow}
\title{Sinkhorn doubly stochastic attention \\rank decay analysis}
\author{Michela Lapenna, Rita Fioresi, Bahman Gharesifard}
\date{Source: arXiv:2604.07925v1 (CC BY 4.0)}
\begin{document}
\maketitle

\noindent\emph{Source and licence.} Sinkhorn doubly stochastic attention \\rank decay analysis. Version arXiv:2604.07925v1 (CC BY 4.0), distributed by arXiv under the Creative Commons Attribution 4.0 International licence, http://creativecommons.org/licenses/by/4.0/, as recorded in the arXiv licence metadata for this exact version and shown on the abstract page https://arxiv.org/abs/2604.07925v1. Full licence notice: \texttt{licenses/arxiv-2604.07925-CC-BY-4.0.txt}.\par\smallskip

\noindent\textit{Editorial scope.} Counted unit: the article's proposition sink\_approx (source0.tex lines 1175--1185). The article's complete original proof of it is delivered with it (source0.tex lines 1187--1250, one \texttt{proof} environment of 64 lines). No mathematics is written, repaired or re-derived: the statement and the proof are the article's own lines, in the article's own notation, and no retained line is substituted or re-flowed.  No further source-local context is needed: every cross-reference of every retained line resolves inside the two delivered spans.  Nothing was omitted from any retained span, so the package carries no omission record.  No retained line contains a citation, so the wrapper carries no bibliography and no bibliography entry.  No retained line contains a figure environment, an image-inclusion command or drawing code, so no image is delivered and the package declares no asset dependency.  Editorial formatting only, in addition: an AMS \texttt{amsart} preamble that restates the article's own declarations and macros used by the retained lines, namely the environments \texttt{proposition} on separate counters, together with the standard packages those lines need.  The retained environments print in this document as Proposition 1 in order of appearance, whereas the article prints the same items with the numbering of its own full document.  The complete native source of the article as distributed in the arXiv e-print for this exact version is bundled unchanged as \texttt{sources/198125/source0.tex}.
\begin{proposition}\label{sink_approx}
Let $F=C \circ R: \R^{n\times n} \lra \R^{n\times n}$ with $R$ row-softmax and $C$ column-softmax operators. Define $Q:\R^{n\times n} \lra \R^{n\times n}$,
$Q(Y) $ = $ Y$
$- \frac{1}{n} \mathbf{1}\mathbf{1}^\top Y
- \frac{1}{n} Y \mathbf{1}\mathbf{1}^\top
+ \frac{1}{n^2} \mathbf{1}\mathbf{1}^\top Y \mathbf{1}\mathbf{1}^\top$.
Then:
\begin{equation}\label{f-op}
F(tX)=\frac{1}{n}\mathbf{1}\mathbf{1}^\top+\frac{t}{n^2}Q(X)+o(t)
\end{equation}
\end{proposition}

\begin{proof}
Let $s: \R^n \lra \R^n$ denote the softmax function. Its Jacobian at $u \in \R^n$ is given by:
$$
J_s(u)=\operatorname{diag}(s(u)) - s(u)s(u)^\top
$$
In particular, evaluating at $u = 0$, we obtain:
\[
J_s(0) = \frac{1}{n} I - \frac{1}{n^2} \mathbf{1}\mathbf{1}^\top= \frac{1}{n} U,
\]
where $U := I - \frac{1}{n} \mathbf{1}\mathbf{1}^\top$.

For a one-parameter perturbation $u(t)=0+t v$ with $v\in\mathbb{R}^n$, we have:
\begin{equation}\label{softeq}
s(tv)=s(0)+tJ_s(0)v+o(t)=\frac{1}{n}\mathbf{1}+\frac{t}{n}Uv+o(t).
\end{equation}

For $Y\in \R^{n\times n}$, define:
$$
r_i(Y):=\begin{pmatrix}y_{i1}\\ \vdots \\ y_{in}\end{pmatrix}\in\mathbb{R}^n,
\qquad
c_j(Y):=\begin{pmatrix}y_{1j}\\ \vdots \\ y_{nj}\end{pmatrix}\in\mathbb{R}^n,
$$

From \ref{softeq}, it follows that:
\[
r_i(R(tX))=s(t\,r_i(X))=\frac{1}{n}\mathbf{1}+\frac{t}{n}U\,r_i(X)+o(t), 
\]
which in matrix form can be written as:
\[
R(tX)=\frac{1}{n}\mathbf{1}\mathbf{1}^\top+\frac{t}{n}XU+o(t). 
\]

We now apply the column-softmax operator to $Z(t):=R(tX)$. We can write: 
\[
Z(t)=Z_0+t\Delta+o(t),
\]
where $Z_0=\frac{1}{n}\mathbf{1}\mathbf{1}^\top$ and $\Delta=\frac{1}{n}XU$.

Since each column of $Z_0$ equals $\frac{1}{n}\mathbf{1}$, and since adding a constant to a column does not change softmax, for each column $j$, we have that: 
\[
c_j(C(Z(t)))=\frac{1}{n}\mathbf{1}+\frac{t}{n}U\,c_j(\Delta)+o(t). 
\]
Written in matrix form, this means that: 
\[
C(Z(t))=\frac{1}{n}\mathbf{1}\mathbf{1}^\top+\frac{t}{n}U\Delta+o(t).
\]
We are now ready to compute:
\[
F(tX)=C(R(tX))=C(Z(t))
=
\frac{1}{n}\mathbf{1}\mathbf{1}^\top+\frac{t}{n}U\Delta+o(t)
=
\frac{1}{n}\mathbf{1}\mathbf{1}^\top+\frac{t}{n}U\Bigl(\frac{1}{n}XU\Bigr)+o(t).
\]
Therefore:
\[
F(tX)
=
\frac{1}{n}\mathbf{1}\mathbf{1}^\top+\frac{t}{n^2}UXU+o(t)
=
\frac{1}{n}\mathbf{1}\mathbf{1}^\top+\frac{t}{n^2}Q(X)+o(t),
\]
where we have used the fact that $Q(X)=UXU$.
\end{proof}

\end{document}
